\documentclass[border=8pt]{standalone}
\usepackage[UTF8]{ctex}
\usepackage{tikz}
\usepackage{amsmath,amssymb}
\usetikzlibrary{arrows.meta,calc}
\definecolor{cBlue}{RGB}{37,99,235}
\definecolor{cOrange}{RGB}{234,88,12}
\definecolor{cGray}{RGB}{100,116,139}
\definecolor{cGreen}{RGB}{22,163,74}
\definecolor{cRed}{RGB}{220,38,38}
\begin{document}
\begin{tikzpicture}[scale=2.6,>=Stealth,line width=0.9pt]
  \draw[cGray,->] (-1.5,0)--(2.1,0) node[below right]{$\mathrm{Re}$};
  \draw[cGray,->] (0,-1.5)--(0,1.7) node[left]{$\mathrm{Im}$};
  \draw[cBlue,line width=1.3pt] (0,0) circle (1);
  \node[cBlue] at (140:1.2) {$U(1)$ （群，弯曲）};
  \fill[cGray] (1,0) circle (0.026);
  \node[cGray,below left=-1pt] at (1,0) {$1$};
  % tangent line at the identity = the Lie algebra
  \draw[cRed,line width=1.2pt] (1,-1.2)--(1,1.3);
  \node[cRed,right] at (1.04,1.18) {$\mathfrak{u}(1)=i\mathbb{R}$ （李代数，平直）};
  % a generator (infinitesimal rotation)
  \draw[cOrange,->,line width=1.3pt] (1,0)--(1,0.8);
  \node[cOrange,right] at (1.04,0.45) {$X=i\,\dot\theta$};
  \fill[cRed] (1,0) circle (0.02);
  \node[align=center] at (0.3,-1.78)
    {\footnotesize\color{cGray} 李代数 $\mathfrak{g}$ $=$ 群在单位元 $1$ 处的切空间；\\[-1pt]
     \footnotesize\color{cGray} 它的元素是“无穷小对称”——这里就是无穷小转动 $i\dot\theta$。};
\end{tikzpicture}
\end{document}
